\section{総括}

本研究では、特異型EEJの発生特性を解明することと、
EE-indexをリアルタイムでモニタリングできるシステムの開発を目的として、
EE-index の開発における精度向上を行った。

次に、特異型 EEJの未発達型と突発型を定量的に分類するために、
自動判定のアルゴリズムを開発した。これにより、精度が向上するとともに、
観測点と期間を任意に指定することが可能になり、
特異型EEJの発生特性を柔軟に調査することが可能になった。

特に南米領域に特化して調査を行った結果、
未発達型では Sq 電流の季節依存性が大きく関与していることが明らかになった。
一方、突発型では Sq 電流の季節依存性と月潮汐効果が組み合わさることによって発生していることが明らかになった。

最後に、EE-index を用いた宇宙天気システムを開発し、
リアルタイムでの EEJ モニタリングを可能とする基盤を構築した。
本システムには、クイックルック機能やデータダウンロード機能など、
EE-index を利用する研究者にとって有用な機能を実装しており、
今後の磁気赤道域研究への貢献が期待される。